\documentclass[../main]{subfiles}
\begin{document}

\chapter{Primer Principio de la Termodinámica}

\section{Teoría del Primer Principio de la Termodinámica}

\subsection{Introducción al Primer Principio de la Termodinámica}

El primer principio de la termodinámica, también conocido como el principio de conservación de la energía, establece que la energía no se crea ni se destruye, solo se transforma de una forma a otra. Este principio es fundamental en la ingeniería y la física, ya que permite entender y analizar los sistemas termodinámicos, desde motores y refrigeradores hasta procesos biológicos y fenómenos atmosféricos.

\subsection{Formulación del Primer Principio de la Termodinámica}

Para un sistema cerrado (un sistema que no intercambia materia con su entorno), el primer principio de la termodinámica se puede expresar matemáticamente como:
\begin{equation}
\Delta U = Q - W
\end{equation}
Donde:
\begin{itemize}
    \item $\Delta U$ es el cambio en la energía interna del sistema.
    \item $Q$ es el calor añadido al sistema.
    \item $W$ es el trabajo realizado por el sistema.
\end{itemize}

El signo de $Q$ y $W$ sigue la convención de signos: el calor añadido al sistema es positivo, y el trabajo realizado por el sistema es positivo.

\subsection{Energía Interna y Capacidades Caloríficas}

La energía interna ($U$) de un sistema es la suma de todas las energías microscópicas (cinéticas y potenciales) de las moléculas que componen el sistema. Un concepto relacionado es la capacidad calorífica, que indica la cantidad de calor necesario para cambiar la temperatura de una cierta cantidad de sustancia en una unidad de temperatura. Para una cantidad de sustancia dada, la capacidad calorífica a volumen constante ($C_V$) y a presión constante ($C_P$) se definen como:
\begin{equation}
C_V = \left( \frac{\partial U}{\partial T} \right)_V
\end{equation}
\begin{equation}
C_P = \left( \frac{\partial H}{\partial T} \right)_P
\end{equation}
Donde $H$ es la entalpía, definida como:
\begin{equation}
H = U + PV
\end{equation}
Donde $P$ es la presión y $V$ es el volumen del sistema.

\subsection{Aplicaciones del Primer Principio de la Termodinámica}

El primer principio se aplica en numerosos campos, incluyendo la ingeniería mecánica y química, la meteorología y la biología. Algunos ejemplos de su aplicación incluyen:

\begin{itemize}
    \item En motores térmicos, el primer principio se usa para analizar la eficiencia del ciclo térmico.
    \item En la refrigeración, se aplica para determinar la cantidad de trabajo necesario para extraer calor de un sistema.
    \item En biología, se utiliza para entender cómo los organismos convierten energía química de los alimentos en trabajo y calor.
\end{itemize}

\actividad{Responder el siguiente cuestionario}
\begin{enumerate}
    \item ¿Qué establece el primer principio de la termodinámica?
    \item ¿Qué se entiende por energía interna de un sistema?
    \item ¿Cómo se define la capacidad calorífica a volumen constante ($C_V$)?
    \item ¿Qué es la entalpía ($H$) y cómo se relaciona con la energía interna ($U$)?
    \item ¿Qué convención de signos se usa para el calor ($Q$) y el trabajo ($W$) en el primer principio de la termodinámica?
    \item ¿Cómo se aplica el primer principio de la termodinámica en un motor térmico?
    \item ¿Qué diferencia hay entre capacidad calorífica a volumen constante ($C_V$) y a presión constante ($C_P$)?
    \item ¿Cómo se puede determinar la eficiencia de un ciclo térmico utilizando el primer principio de la termodinámica?
    \item ¿En qué consiste un sistema cerrado en termodinámica?
    \item ¿Qué papel juega el primer principio de la termodinámica en la refrigeración?

\end{enumerate}

\actividad{Resolver los siguientes problemas}
\begin{enumerate}
    \item Un gas ideal se encuentra en un recipiente cerrado y recibe 500 J de calor mientras realiza 200 J de trabajo. ¿Cuál es el cambio en la energía interna del gas?
    \item Un sistema absorbe 1500 J de calor y su energía interna aumenta en 1000 J. ¿Cuánto trabajo realiza el sistema?
    \item Se añade 800 J de calor a un gas y se realiza un trabajo de 600 J sobre el gas. ¿Cuál es el cambio en la energía interna del gas?
    \item Un sistema realiza 400 J de trabajo y su energía interna disminuye en 700 J. ¿Cuánto calor ha intercambiado el sistema con su entorno?
    \item Un gas en un recipiente rígido recibe 1000 J de calor. ¿Cuál es el cambio en la energía interna del gas?
    \item Un gas ideal en un recipiente cerrado y rígido tiene una capacidad calorífica a volumen constante ($C_V$) de 20 J/K. Si se añade 300 J de calor al gas, ¿cuál es el cambio en su temperatura?
    \item Un sistema cerrado absorbe 1200 J de calor mientras realiza 400 J de trabajo. ¿Cuál es el cambio en la energía interna del sistema?
    \item Se realiza un trabajo de 500 J sobre un gas ideal y su energía interna aumenta en 300 J. ¿Cuánto calor ha intercambiado el gas con su entorno?
    \item Un gas realiza 250 J de trabajo y su energía interna aumenta en 500 J. ¿Cuánto calor ha recibido el gas?
    \item Un sistema cerrado pierde 800 J de energía interna mientras realiza 200 J de trabajo. ¿Cuánto calor ha intercambiado con su entorno?
\end{enumerate}

\end{document}
